\SourceLeaf{482}wegen oftmaliger Veränderungen an den Lägern und Pfannen an den verschiedenen Tagen nicht gleiche Werthe gehabt haben. Meistens sind die Beobachtungen in den Vormittagsstunden zwischen 8 und 11 Uhr angestellt; am 16.\ 22.\ 25.\ Juni und 17.\ 20.\ Juli aber Nachmittags zwischen 4 und 6 Uhr.

Die einzelnen Columnen geben an: das Zeichen des Nordpolendes der Nadel, die Werthe von $f$ und $g$ oder von $f'$ und $g'$, je nachdem $B$ oder $A$ der Nordpol gewesen, und die Dauer der horizontalen Schwingung.

\Needspace{16\baselineskip}
\begin{center}
\textit{Beobachtungen mit Nadel} 1.

\begin{tabular}{rr|c|rrr|rrr|l}
Mai&20&$B$&$67^0$&$11'$&$0''$&$67^0$&$58'$&$46''$&$5''87152$\\
&21&$A$&&$57$&$1$&&$35$&$14$&$5.81508$\\
&22&$A$&&$56$&$29$&&$36$&$45$&$5.82044$\\
&24&$B$&&$16$&$45$&&$45$&$48$&$5.81557$\\
&31&$B$&&$18$&$1$&&$49$&$41$&$5.82075$\\
Juni&2&$A$&&$53$&$55$&&$33$&$9$&$5.85778$\\
&4&$A$&&$56$&$38$&&$32$&$10$&$5.86442$\\
&5&$B$&&$24$&$13$&&$46$&$44$&$5.83615$\\
Juli&6&$A$&&$59$&$41$&&$35$&$21$&$5.83716$\\
&7&$A$&&$58$&$7$&&$37$&$51$&$5.83818$\\
&8&$B$&&$20$&$8$&&$44$&$47$&$5.89602$\\
&9&$B$&&$20$&$43$&&$44$&$25$&$5.90035$
\end{tabular}
\end{center}

\Needspace{16\baselineskip}
\begin{center}
\textit{Beobachtungen mit Nadel} 2.

\begin{tabular}{rr|c|rrr|rrr|l}
Mai&20&$A$&$67^0$&$40'$&$57''$&$67^0$&$20'$&$37''$&$5''72416$\\
&21&$A$&&$41$&$8$&&$21$&$5$&$5.72453$\\
&22&$B$&&$43$&$28$&&$50$&$45$&$5.65355$\\
&24&$B$&&$41$&$43$&&$54$&$32$&$5.66875$\\
&31&$A$&&$43$&$34$&&$18$&$29$&$5.67439$\\
Juni&2&$A$&&$41$&$46$&&$18$&$12$&$5.67665$\\
&4&$B$&&$42$&$42$&&$46$&$57$&$5.68010$\\
&5&$B$&&$44$&$53$&&$50$&$24$&$5.68890$\\
Juli&17&$B$&&$45$&$20$&&$50$&$17$&$5.70183$\\
&18&$A$&&$40$&$26$&&$22$&$50$&$5.68692$\\
&19&$A$&&$40$&$21$&&$22$&$10$&$5.69677$\\
&20&$B$&&$40$&$40$&&$54$&$19$&$5.66585$
\end{tabular}
\end{center}

\SourceLeaf{483}\Needspace{16\baselineskip}
\begin{center}
\textit{Beobachtungen mit Nadel} 3.

\begin{tabular}{rr|c|rrr|rrr|l}
Juni&8&$B$&$67^0$&$47'$&$58''$&$67^0$&$48'$&$52''$&$6''17149$\\
&9&$B$&&$40$&$55$&&$42$&$28$&$6.18077$\\
&11&$A$&&$30$&$58$&&$32$&$35$&$6.18080$\\
&16&$B$&&$40$&$0$&&$42$&$40$&$6.17046$\\
&18&$B$&&$43$&$13$&&$47$&$40$&$6.18005$\\
&22&$A$&&$27$&$33$&&$39$&$19$&$6.16591$\\
&23&$A$&&$29$&$46$&&$41$&$8$&$6.16948$\\
&25&$A$&&$29$&$3$&&$41$&$7$&$6.17663$\\
Juli&6&$A$&&$32$&$38$&&$40$&$37$&$6.18305$\\
&7&$B$&&$45$&$56$&&$42$&$12$&$6.17982$\\
&8&$B$&&$46$&$59$&&$43$&$37$&$6.18339$\\
&9&$A$&&$30$&$42$&&$39$&$42$&$6.23905$
\end{tabular}
\end{center}

\Needspace{16\baselineskip}
\begin{center}
\textit{Beobachtungen mit Nadel} 4.

\begin{tabular}{rr|c|rrr|rrr|l}
Juni&8&$A$&$67^0$&$45'$&$9''$&$67^0$&$27'$&$3''$&$5''96200$\\
&9&$B$&&$22$&$56$&$68$&$8$&$28$&$5.91653$\\
&11&$B$&&$23$&$16$&&$7$&$48$&$5.94665$\\
&16&$A$&&$49$&$54$&$67$&$12$&$8$&$6.01785$\\
&18&$B$&&$27$&$48$&$68$&$8$&$45$&$5.93204$\\
&22&$B$&&$26$&$46$&&$3$&$56$&$5.94065$\\
&23&$A$&&$50$&$19$&$67$&$15$&$37$&$5.93939$\\
&25&$A$&&$50$&$4$&&$15$&$22$&$5.94731$\\
Juli&17&$A$&&$50$&$13$&&$15$&$43$&$5.96850$\\
&18&$A$&&$49$&$57$&&$14$&$48$&$5.96931$\\
&19&$B$&&$22$&$43$&$68$&$9$&$18$&$5.92673$\\
&20&$B$&&$22$&$41$&&$10$&$19$&$5.92783$
\end{tabular}
\end{center}

\Article{24}
Bei der Berechnung dieser Beobachtungen werde ich anstatt der oben (Art.\ 21.\ 22) gebrauchten $t,u$ etwas modificirte Hülfsgrössen einführen. Wenn man für eine der Nadeln die Dauer einer horizontalen Schwingung mit $n$, die Summe der Trägheitsmomente der Nadel und des Bügels in Beziehung auf die bei diesen Schwingungen verticale Drehungsachse mit $k$, und die Länge des einfachen Secundenpendels mit $l$ bezeichnet. so ist bekanntlich
\[
lmnn\cos i=k
\]
\SourceLeaf{484}Man wähle eine Normalschwingungsdauer $N$ und eine Normalinclination, die zwischen den vorgekommenen Werthen von $n$ und $i$ ungefähr das Mittel halten, und bezeichne den entsprechenden Werth von $m$ mit $M$, so dass
\[
lMNN\cos I=k
\]
wird. Endlich sei
\[
\begin{aligned}
x&=\frac{q\cos Q.\cos I.206265''}{M}\\
y&=\frac{q\sin Q.\sin I.206265''}{M}
\end{aligned}
\]
welche Grössen also für alle Beobachtungen mit dieser Nadel constant sind. Die Gleichungen werden dann
\[
\begin{aligned}
i&=f+c-\frac{nn\cos f}{NN\cos I}.\frac{\cos i}{\cos I}.x+\frac{nn\sin f}{NN\sin I}.\frac{\cos i}{\cos I}.y\\
i&=g-c-\frac{nn\cos g}{NN\cos I}.\frac{\cos i}{\cos I}.x-\frac{nn\sin g}{NN\sin I}.\frac{\cos i}{\cos I}.y
\end{aligned}
\]
wenn $B$ der Nordpol ist; für den Fall, wo $A$ der Nordpol ist, hat man nur den $x$ und $y$ enthaltenden Gliedern die entgegengesetzten Zeichen zu geben.

Diese Form hat den Vortheil, dass die Coëfficienten von $x$ und $y$ immer wenig von der Einheit verschieden sind, und in der That kann man bei so geringer Excentricität des Schwerpunkts, wie die vier in Rede stehenden Nadeln haben, und bei so mässigen Schwankungen von $n$, anstatt jener Coëfficienten füglich die Einheit annehmen, welches ich die abgekürzte Rechnung nenne. Indessen habe ich mir doch die Mühe gegeben, die $192$ Coëfficienten genauer zu berechnen und nur den Factor $\frac{\cos i}{\cos I}$ weggelassen, wenn auch der Nutzen davon hauptsächlich nur darin besteht, die Zulässigkeit der abgekürzten Rechnung desto anschaulicher zu machen. Fortan sollen die nichtaccentuirten Buchstaben $N,x,y$ sich auf die Nadel $1$ beziehen, und die Werthe für die drei andern Nadeln der Reihe nach durch einen, zwei und drei Accente unterschieden werden. Gewählt sind für gegenwärtige Rechnung die Werthe
\[
\begin{aligned}
I&=67^0\ 40'\ 0''\\
N&=5''847785\\
N'&=5.686867\\
N''&=6.181742\\
N'''&=5.949567
\end{aligned}
\]
\expandafter\def\csname HistoricalNoteText485\endcsname{Es bedarf keiner Erinnerung, dass ein solcher Fehler, der, wenn er überhaupt reell ist, nur einer Abweichung der Zapfen von der cylindrischen Gestalt zugeschrieben werden kann, nur in sofern constant ist, als immer dieselben Stellen der Zapfen zum Aufliegen kommen, also bei einer ganz andern Inclination auch einen ganz verschiedenen Werth haben könnte.}
\SourceLeaf{485}Die Rechnungen selbst werde ich, um den Raum zu schonen, hier nicht in extenso aufnehmen, sondern nur so viel davon mittheilen, als nöthig ist, um dem Gange im Allgemeinen folgen zu können. Übrigens sind die von der Einheit am meisten abweichenden Werthe der Coëfficienten $0.96895$ und $1.04324$, welche am 9. und 16. Juni bei Nadel 4 vorkommen.

\Article{25}
Aus den beiden Gleichungen, welche die Beobachtungen mit einer Nadel an jedem Tage liefern, bilden sich, indem man sowohl ihre Summe als ihre Differenz halbirt, zwei andere, die mit I und II bezeichnet werden mögen. Es entstehen also $48$ Gleichungen I, und eben so viele II, von denen ich die ersten als Probe hersetze. Die ursprünglichen Gleichungen aus den Beobachtungen vom 20. Mai mit Nadel 1 sind
\[
\begin{aligned}
i&=67^0\ 11'\ 0''+c-1{,}02880x+1{,}00460y\\
i&=67\ 58\ 46-c-0{,}99473x-1{,}01038y
\end{aligned}
\]
woraus die abgeleiteten entstehen
\[
\begin{aligned}
i&=67^0\ 34'\ 53''-1{,}01176x-0{,}00289y\quad.\ .\ .\ .\ .\ .\quad\text{(I)}\\
c&=+1433''+0{,}01703x-1{,}00749y\quad.\ .\ .\ .\ .\ .\quad\text{(II)}
\end{aligned}
\]
Um die im 8. Art. angedeutete Prüfung anstellen zu können, habe ich aber den Gleichungen I noch ein Glied beigefügt, indem ich $i+e$ anstatt $i$ schreibe, so dass $e$ den etwanigen constanten\HistoricalNote{485} Fehler der Nadel 1 ausdrückt; bei den Nadeln 2, 3, 4 soll der präsumtive constante Fehler mit $e',e'',e'''$ bezeichnet werden.

Auf diese Weise schliessen also die $48$ Gleichungen I zusammen $36$ unbekannte Grössen ein, nemlich die Inclinationen an den 24 Beobachtungstagen, und die $12$ Grössen $x,y,e,x',y',e',x''$ u. s. w. Es muss aber zuvörderst bemerkt werden, dass die Glieder, welche $y,y',y'',y'''$ enthalten, alle nur sehr kleine Coëfficienten haben, und in der abgekürzten Rechnung ganz fehlen: der grösste dieser $48$ Coëfficienten ist eben $0.00289$ in der obigen Probegleichung. Will
\expandafter\def\csname HistoricalNoteText486\endcsname{Die vom 23. September ist die, welche oben Art.\ $9$–$22$ als Beispiel gedient hat; die beiden andern werden unten Art.\ $30$ angeführt.}
\SourceLeaf{486}man aber einmal den geringen nur wenige Secunden betragenden Einfluss berücksichtigen, so muss man zuvor die Werthe dieser $y,y',y'',y'''$ anderswoher abgeleitet haben, wo aber jedenfalls grob genäherte Werthe zu diesem Zweck schon zureichend sind.

\Article{26}
Zu dieser Ableitung stehen uns nun nur die Gleichungen II zu Gebote. Allein wenn man erwägt, dass in den $12$ Gleichungen dieser Abtheilung, welche sich auf Eine Nadel beziehen, der Buchstab $c$ ungleiche Werthe repräsentirt, indem bei jedem Umstreichen der Werth verändert werden kann, so erkennt man leicht, dass es unmöglich ist, diese $c$ aus den Gleichungen zu eliminiren, und dass man also \textit{gezwungen} ist, eine etwas precäre Hypothese zu Hülfe zu nehmen. Die meinige besteht in Folgendem. Da, bei allen bedeutenden Schwankungen von $c$, doch unter Anwendung eines immer gleichen Streichverfahrens ein Mittelwerth von $c$ sich herausstellen wird, so nehme ich an, dass der Mittelwerth für die eine Lage der Pole derselbe ist wie für die andere. Freilich wird nur eine sehr unvollkommene Compensation zu erwarten sein, wenn nur eine geringe Anzahl von Umstreichungen Statt gefunden hat, und der auf diese Weise abgeleitete Werth von $y$ wird also wenig Sicherheit haben; allein dieser Unsicherheit ist gar nicht auszuweichen, wenn man nicht die Werthe von $c$ durch einen besondern Apparat ausmittelt (S.\ oben Art.\ 17). Zur Benutzung jenes Princips wird man also bei jeder Nadel zuerst die Gleichungen II, welche sich auf $B$ Nord beziehen, von denen trennen, wo $A$ Nord war; dann die erstern und die letztern in so viele Gruppen zerlegen, als veränderte magnetische Zustände Statt gefunden haben; aus den zu derselben Gruppe gehörenden Gleichungen (in sofern mehrere in Eine Gruppe kommen) das Mittel, und aus diesen partiellen Mitteln wieder das Mittel nehmen; indem man dann die so hervorgehenden Mittelwerthe einander gleich setzt, erhält man die Gleichung, durch welche $y$ bestimmt wird. Zur Erläuterung setze ich die \textit{abgekürzte} Rechnung für Nadel 1 her, bei welcher ich \textit{zu diesem Zwecke} obigen $12$ Beobachtungen auch noch drei andere\HistoricalNote{486} vom 1. August, 7. August, 23. September benutzt habe. Während des ganzen Zeit\SourceLeaf{487}raumes war die Nadel neunmal umgestrichen, so dass zehn verschiedene Zustände Statt gefunden haben, wovon fünf auf jede Lage der Pole kommen.

\Needspace{16\baselineskip}
\begin{center}\textit{Nadel} 1, $B$ \textit{Nord}\end{center}
\[
\begin{array}{rr|l}
&&\qquad c+y=\\\hline
\text{Mai}&20&.\ .\ .\quad+1433''\\
&\begin{array}{r}24\\31\end{array}&\left.\begin{array}{r}+871''\\+950\end{array}\right\}+910\\
\text{Juni}&5&.\ .\ .\quad+675\\
\begin{array}{r}\text{Juli}\\\phantom{\text{Juli}}\\\text{Aug.}\end{array}&\begin{array}{r}8\\9\\1\end{array}&\left.\begin{array}{r}+739\\+711\\+720\end{array}\right\}+723\\
\begin{array}{r}\phantom{\text{Aug.}}\\\text{Sept.}\end{array}&\begin{array}{r}7\\23\end{array}&\left.\begin{array}{r}+584\\+528\end{array}\right\}+556\\\cline{3-3}
\multicolumn{2}{r}{\text{Mittel}}&c+y=+859''
\end{array}
\]
\Needspace{16\baselineskip}
\begin{center}\textit{Nadel} 1, $A$ \textit{Nord}\end{center}
\[
\begin{array}{rr|l}
&&\\\hline
\text{Mai}&\begin{array}{r}21\\22\end{array}&\left.\begin{array}{r}-653''\\-592\end{array}\right\}-623''\\
\text{Juni}&\begin{array}{r}2\\4\end{array}&\left.\begin{array}{r}-623\\-734\end{array}\right\}-678\\
\text{Juli}&\begin{array}{r}6\\7\end{array}&\left.\begin{array}{r}-730\\-608\end{array}\right\}-669\\
\text{Aug.}&\begin{array}{r}1\\7\end{array}&\left.\begin{array}{r}-720\\-785\end{array}\right\}-752\\
\text{Sept.}&23&.\ .\ .\quad-680\\\cline{3-3}
\multicolumn{2}{r}{\text{Mittel}}&c-y=-680
\end{array}
\]
woraus also $y=+769''$ folgt. Die nicht abgekürzte Rechnung ergab
\[
\begin{aligned}
\text{für }B\text{ Nord},\quad c&=+859+0.00102x-1.00290y\\
\text{für }A\text{ Nord},\quad c&=-680+0.00082x+0.99915y
\end{aligned}
\]
woraus
\[
y=+769''+0.00097x
\]
folgt. Auf gleiche Weise findet sich für die drei andern Nadeln
\SourceLeaf{488}\[
\begin{aligned}
y'&=+456''-0.00192x'\\
y''&=-101+0.00134x''\\
y'''&=+1107+0.00224x'''
\end{aligned}
\]
Die Schwankungen in den Werthen von $c$ gehen bei der Nadel 1 auf $14\frac12$ Minuten, bei den Nadeln 2 und 3 auf $4\frac12$ Minuten, bei der Nadel 4 auf $10$ Minuten. Damit man übrigens dem Umstande, dass gerade an dem ersten Beobachtungstage der am meisten abweichende Werth bei der Nadel 1 vorkommt, nicht eine besondere Wichtigkeit beilege, will ich noch bemerken, dass sowohl an dieser, wie an den übrigen Nadeln die Pole vor den hier mitgetheilten Beobachtungen schon oft und immer mit derselben Sorgfalt und denselben Streichmitteln umgekehrt gewesen waren.

\Article{27}
Nachdem die Werthe von $y,y',y'',y'''$ in den Gleichungen I substituirt sind, bleiben in denselben noch $32$ unbekannte Grössen, und wenn man dann immer die beiden Gleichungen, welche für die Beobachtungen eines und desselben Tages gelten, von einander abzieht, so bilden sich $24$ neue Gleichungen, welche nur die acht unbekannten Grössen $x,x',x'',x''',e,e',e'',e'''$ enthalten. Die vier letzten kommen aber nur in den Differenzen von je zweien vor, so dass man, wenn man
\[
\begin{aligned}
e'-e&=d'\\
e''-e&=d''\\
e'''-e&=d'''
\end{aligned}
\]
setzt, nur sieben unbekannte Grössen behält. Die Coëfficienten von $d',d'',d'''$ sind darin alle $+1$ oder $-1$, und die Coëfficienten von $x,x',x'',x'''$ alle von $+1$ oder $-1$ sehr wenig verschieden. Zur Bestimmung der Werthe der sieben unbekannten Grössen vermittelst der Methode der kleinsten Quadrate wird man, Behuf der Bildung der auf $x,x',x'',x'''$ sich beziehenden Normalgleichungen, die Multiplication mit den respectiven Coëfficienten ohne Bedenken unterlassen können, so dass zur Bildung sämmtlicher sieben Normalgleichungen nichts als einfache Addition erforderlich ist. Auf diese Art haben sich folgende Normalgleichungen ergeben:
\SourceLeaf{489}\[
\begin{aligned}
0&=+4804''+12.00266x-0.00708x'+0.01900x''\\
0&=-5806+0.01559x+12.01005x'-0.00072x'''\\
0&=-3228+0.00145x+12.00544x''+0.04561x'''\\
0&=-5267+0.01786x'-0.00489x''+12.00343x'''\\
0&=-297+0.02717x+0.11088x'-0.04723x'''\\
 &\qquad-12d'+4d'''\\
0&=-241+0.06326x+0.05839x''-0.08085x'''\\
 &\qquad-12d''+8d'''\\
0&=+254-0.02682x'-0.02676x''+0.12808x'''\\
 &\qquad+4d'+8d''-12d'''
\end{aligned}
\]
und hieraus die Werthe
\[
\begin{aligned}
x&=-400''\\
x'&=+484\\
x''&=+267\\
x'''&=+438\\
d'&=-22\\
d''&=-23\\
d'''&=+1
\end{aligned}
\]
Anstatt der drei letzten kann man auch, indem man
\[
\tfrac14(e+e'+e''+e''')=\varepsilon
\]
setzt, schreiben
\[
\begin{aligned}
e&=+11''+\varepsilon\\
e'&=-11+\varepsilon\\
e''&=-12+\varepsilon\\
e'''&=+12+\varepsilon
\end{aligned}
\]
wo der gemeinschaftliche Theil $\varepsilon$ offenbar aus den zu Gebote stehenden Daten nicht bestimmbar ist. Die Substitution der gefundenen Werthe von $x,e,x',e'$ u. s. w. in den (von $y,y'$ u. s. w. bereits befreieten) Gleichungen I gibt uns nun, unter Weglassung von $\varepsilon$, folgende $48$ Inclinationen:
\SourceLeaf{490}\Needspace{29\baselineskip}
\begin{center}
\begin{tabular}{rr|c|rrr||c|rrr}
&&Nadel&&&&Nadel&&&\\\hline
Mai&20&$1$&$67^0$&$41'$&$25''$&$2$&$67^0$&$39'$&$12''$\\
&21&&&$39$&$21$&&&$39$&$31$\\
&22&&&$39$&$51$&&&$39$&$22$\\
&24&&&$37$&$43$&&&$40$&$21$\\
&31&&&$40$&$17$&&&$39$&$17$\\
Juni&2&&&$36$&$39$&&&$38$&$16$\\
&4&&&$37$&$31$&&&$37$&$0$\\
&5&&&$41$&$56$&&&$39$&$48$\\
&8&$3$&&$44$&$12$&$4$&&$43$&$14$\\
&9&&&$37$&$27$&&&$38$&$15$\\
&11&&&$36$&$27$&&&$38$&$1$\\
&16&&&$37$&$6$&&&$38$&$17$\\
&18&&&$41$&$12$&&&$40$&$48$\\
&22&&&$38$&$5$&&&$37$&$51$\\
&23&&&$40$&$6$&&&$40$&$2$\\
&25&&&$39$&$45$&&&$39$&$49$\\
Juli&6&$1$&&$40$&$42$&$3$&&$41$&$17$\\
&7&&&$41$&$11$&&&$39$&$49$\\
&8&&&$39$&$5$&&&$41$&$3$\\
&9&&&$39$&$12$&&&$39$&$57$\\
&17&$2$&&$39$&$55$&$4$&&$40$&$7$\\
&18&&&$39$&$56$&&&$39$&$31$\\
&19&&&$39$&$35$&&&$38$&$32$\\
&20&&&$39$&$43$&&&$39$&$1$
\end{tabular}
\end{center}

\Article{28}
Die Ungleichheiten zwischen den beiden Bestimmungen der Inclination an jedem Tage werden uns nun den Maassstab für die Unsicherheit der Beobachtungen selbst geben müssen. Die grösste Ungleichheit (am 24. Mai) beträgt $2'\ 38''$, und die Summe der Quadrate aller $24$ Unterschiede, die Secunde als Einheit angenommen, ist $124389$. Aus den Principien der Wahrscheinlichkeitsrechnung ist leicht abzuleiten, dass wenn wir den Beobachtungen mit den einzelnen vier Nadeln gleiche Zuverlässigkeit beilegen (von welcher Voraussetzung abzugehen keine Gründe vorhanden sind), die mittlere Unsicherheit eines aus den Beobachtungen gefundenen und unsern Rechnungen untergelegten Werthes von $\tfrac12(f+g)$ oder $\tfrac12(f'+g')$, so weit sich darüber nach unsern Zahlen urtheilen lässt,
\[
=\BareRoot{\frac{124389}{34}}=60''5
\]
\SourceLeaf{491}gesetzt werden muss, insofern nemlich nur von den zufälligen oder regellosen Beobachtungsfehlern die Rede ist. Das Mittel aus zwei solchen auf von einander unabhängige Beobachtungen gegründeten Zahlen wird folglich mit der mittlern Unzuverlässigkeit
\[
=\BareRoot{\frac{124389}{68}}=42''8
\]
behaftet sein, und dies kann auch wie der mittlere Fehler einer auf die gewöhnliche Art (d.\ i.\ mit Einer Nadel aber in \textit{beiden} Lagen der Pole) bestimmten Inclination betrachtet werden, insofern die kleine zu $\tfrac14(f+g+f'+g')$ hinzukommende Correction entweder für ganz unmerklich gilt, oder auf sonst schon feststehende Bestimmung von $u$ oder $y$ gegründet werden kann (vergl.\ Art.\ 21). Es versteht sich von selbst, dass diese Fehlerschätzung zunächst nur für dieses Instrument und für solche Beobachtungen gilt, die unter ganz ähnlichen Umständen gemacht sind, wie die zum Grunde liegenden. Bei einer geringern Anzahl von Einstellungen, als acht in jeder Combination, würde die Zuverlässigkeit geringer sein, obwohl ich nicht behaupten möchte, dass der mittlere Fehler des Endresultats genau im verkehrten Verhältnisse der Quadratwurzel aus der Zahl der mit Pfannen vervielfältigten Einstellungen stehe. Von der andern Seite darf ich nicht unbemerkt lassen, dass während der ganzen Dauer obiger Beobachtungen die Läger nicht so vollkommen berichtigt werden konnten, wie ich wünschte, und nachher durch Anwendung des oben (Art.\ 5) erwähnten Apparats wirklich erreichte: die aus einer unvollkommenen Lagerberichtigung möglicher Weise entspringende Vergrösserung der Beobachtungsfehler (wobei an einen Einfluss von constanter Grösse um so weniger zu denken ist, weil sehr oft an den Lägern Veränderungen gemacht wurden) ist demnach in obiger Zahl schon mit begriffen, und ich habe daher Grund zu erwarten, dass künftige Beobachtungen mit demselben Instrument eher noch kleinere Fehler zeigen werden.

Eine besondere Untersuchung, deren Einzelnes ich hier übergehe, hat übrigens ergeben, dass die mittlere Unsicherheit der im vorhergehenden Art.\ angegebenen $48$ Inclinationen nicht viel von der mittlern Unsicherheit der $\tfrac12(f+g)$ verschieden ist, und dass den im 30. Art. zusammenzustellenden Mitteln aus jedem zusammengehörenden Paare nahe das doppelte Gewicht, also der mittlere Fehler $42''8$, beigelegt werden muss.
